\section{Effective uniform saturation and a growing-index zero regime}
\label{efflerch:sec:main}

The fixed-index concentration theorem in the manuscript gives an
existential derivative-order threshold and confines its expansion to a
compact interval of the scale $a/k$.  The estimates below remove both
restrictions for a leading-order location theorem.  They give an explicit,
deliberately conservative threshold in terms of the Laurent index alone.
The proof treats the smallest and largest zero on the same footing.

Retain the reciprocal-gamma polynomials and the positive Lerch weight
\[
 \frac{e^{xz}}{\Gamma(1+z)}
   =\sum_{n\geq0}Q_n(x)\frac{z^n}{n!},
 \qquad W_\rho(t)=\frac1{1-\rho e^{-t}},\quad 0\leq\rho\leq1.
\]
For $k\geq1$, write
\begin{equation}\label{efflerch:eq:F}
 F_{n,k}^\rho(a)=(-1)^{n+k}D_{n,k}^\rho(a)
 =\int_0^\infty t^k e^{-at}W_\rho(t)Q_n(\log t)\,dt.
\end{equation}
Here $D_{n,k}^\rho=\partial_a^k\ell_n(\rho,a)$ for $\rho<1$
and $D_{n,k}^1=\gamma_n^{(k)}(a)$, with
\[
 \ell_n(\rho,a)=\sum_{m\geq0}
       \frac{\rho^m\log^n(a+m)}{a+m}
 =(-1)^n\left.\partial_s^n\Phi(\rho,s,a)\right|_{s=1}.
\]
The integral representation and the
global bound of at most $n$ positive zeros, counted with multiplicity,
are existing results of the manuscript.  We use them explicitly;
the new argument concerns quantitative separation and uniform sign
stability.

Order the roots of $Q_n$ decreasingly as
$x_{n,1}>\cdots>x_{n,n}$ and put $c_{n,j}=e^{-x_{n,j}}$.

\begin{theorem}[Explicit saturation and simultaneous location bounds]
\label{efflerch:thm:explicit}
For every integer $n\geq2$, define
\begin{equation}\label{efflerch:eq:Kn}
 \delta_n=(3n)^{-(n-2)},\qquad
 K_n=576n^2(3n)^{2n-4}=(24n/\delta_n)^2.
\end{equation}
For every integer $k\geq K_n$ and every $0\leq\rho\leq1$,
$D_{n,k}^\rho$ has exactly $n$ positive zeros.  All are simple.
If they are ordered increasingly as $a_{n,1,k}^\rho<\cdots<a_{n,n,k}^\rho$,
then
\begin{equation}\label{efflerch:eq:location}
 \left|\log\frac{a_{n,j,k}^\rho}{k c_{n,j}}\right|
   <\frac{8n}{\sqrt{k}},\qquad 1\leq j\leq n.
\end{equation}
Thus the relative bounds
\[
 e^{-8n/\sqrt{k}}k c_{n,j}
  <a_{n,j,k}^\rho<e^{8n/\sqrt{k}}k c_{n,j}
\]
hold simultaneously for every branch and every deformation parameter.
No compactness assumption on the collection of slopes $c_{n,j}$ is made.
\end{theorem}

The threshold is an upper bound, not a proposed sharp value.  In particular,
the existing index-two theorem gives saturation already at $k=1$, whereas
\eqref{efflerch:eq:Kn} gives $K_2=2304$.  The value of the bound is its
uniformity and explicit dependence on $n$.

\subsection{An explicit separation bound for the Appell roots}

For a real-rooted polynomial, its \emph{mesh} is the minimum gap between
consecutive roots, counted with multiplicity.  A multiple root gives mesh
zero.  During a temporary repeated-root stage, we separately refer to
the minimum gap between its \emph{distinct} roots.

We need the classical fact that $p\mapsto p+cp'$ does not decrease
mesh.  This is the theorem of Riesz recorded as Theorem~1 of
Br\"and\'en--Krasikov--Shapiro~\cite{efflerch:BKS}.  The following proof
also explains why it applies to the infinite reciprocal-gamma product.

\begin{lemma}[Preservation of mesh]
\label{efflerch:lem:mesh-preserve}
If $p$ is real-rooted and $c\in\mathbb R$, then
\[
 \operatorname{mesh}(p+cp')\geq\operatorname{mesh}(p).
\]
\end{lemma}
\begin{proof}
The operator $T=1+c\partial_x$ preserves real-rootedness: away from
the distinct roots $r_i$ of multiplicities $m_i$, its new roots solve
$\sum_i m_i/(x-r_i)=-1/c$.  This sum is strictly decreasing on each
complementary interval.  Its inherited roots have multiplicity $m_i-1$.
The assertion is immediate when $c=0$ or $p$ has mesh zero.

Suppose the mesh of $p$ is $d>0$.  For any $0<h<d$, the roots of $p(x)$
and $p(x-h)$ interlace.  The elementary real-rooted-pencil criterion
says that two real-rooted polynomials of the same degree interlace
if and only if every real linear combination is real-rooted or zero.
One proof divides one polynomial by the other: the interlacing case
has residues of one sign, which makes the rational quotient strictly
monotone between poles.  Conversely, if every pencil member is
real-rooted, that quotient cannot take a real value in the upper
half-plane; its imaginary part has constant sign there, forcing the
residues, and hence the interlacing order.  Common roots follow by
cancellation or a limiting argument.

Apply $T$ to this pencil.  Since $T$ preserves real-rootedness and
commutes with translation, every real combination of $Tp(x)$ and
$Tp(x-h)$ is real-rooted or zero.  The criterion gives interlacing
of this shifted pair, which is equivalent to
$\operatorname{mesh}(Tp)\geq h$.  Let $h\uparrow d$.
\end{proof}

\begin{lemma}[Quantitative separation while removing multiplicities]
\label{efflerch:lem:split}
Let $p$ be real-rooted of degree $n\geq2$, with at least two distinct
roots separated by at least $d$, where $0<d\leq1$.  If $c\geq1/n$,
the distinct roots of $p+cp'$ are separated by at least $d/(3n)$.
\end{lemma}
\begin{proof}
Write the distinct roots as $r_1<\cdots<r_q$, with multiplicities $m_i$.
There is one new root $s_i$ in each gap $(r_i,r_{i+1})$, one new root
$s_0<r_1$, and the inherited roots of reduced multiplicity.  In an
interior gap put $u=s_i-r_i$ and $v=r_{i+1}-s_i$.

If $u<d/2$, every root to the right is at distance at least $d/2$.
In the equation $p'(s_i)/p(s_i)=-1/c$, the positive part includes
$m_i/u$, and the magnitude of the negative part is at most $2n/d$.
Consequently $u\geq d/(2n)$.  This same lower bound is immediate
if $u\geq d/2$.
If $v<d/2$, the positive part is at most $2n/d$, so the negative
part, which includes $m_{i+1}/v$, is at most $2n/d+1/c$.
Thus in every case
\[
 u\geq\frac{d}{2n},\qquad
 v\geq\frac{d}{2n+d/c}\geq\frac{d}{3n}.
\]
For the exterior root,
$\sum_i m_i/(r_i-s_0)=1/c$ implies $r_1-s_0\geq cm_1\geq1/n$.
Every possible pair of consecutive distinct new or inherited roots
is therefore separated by at least $d/(3n)$.
\end{proof}

\begin{proposition}[A fully explicit Appell mesh bound]
\label{efflerch:prop:mesh}
For every $n\geq2$,
\begin{equation}\label{efflerch:eq:mesh}
 \operatorname{mesh}(Q_n)\geq\delta_n=(3n)^{-(n-2)}.
\end{equation}
\end{proposition}
\begin{proof}
Start with $x^n$ and apply $1+\partial_x$.  The result
$x^{n-1}(x+n)$ has two distinct roots at distance $n\geq1$.
Apply next $1+\partial_x/j$, for $j=2,\ldots,n-1$.
Their coefficients $1/j$ are all at least $1/n$.
Lemma~\ref{efflerch:lem:split}, used $n-2$ times starting with $d=1$,
shows that the distinct-root gaps of
\[
 B_n(x)=\prod_{j=1}^{n-1}(1+\partial_x/j)x^n
\]
are at least $(3n)^{-(n-2)}$.  At each step inherited multiplicities
drop by one and all new roots are simple.  After these $n-1$ steps,
every root is simple, so this is a bound for the ordinary mesh of $B_n$.

For $N\geq n-1$ the reciprocal-gamma finite-product polynomial is
\[
 Q_{n,N}(x)=
  \left(\prod_{j=1}^N(1+\partial_x/j)\right)(x+\gamma-H_N)^n.
\]
Translations preserve mesh.  The remaining factors, with $j\geq n$,
do not decrease it by Lemma~\ref{efflerch:lem:mesh-preserve}.
Hence every $Q_{n,N}$ has mesh at least $\delta_n$.
The Weierstrass product
\[
 \Gamma(1+z)^{-1}=e^{\gamma z}
   \prod_{j\geq1}(1+z/j)e^{-z/j}
\]
converges locally uniformly.  Thus $Q_{n,N}\to Q_n$ coefficientwise.
Continuity of the ordered real roots of monic polynomials of fixed
degree preserves the gap inequalities in the limit.  This proves
\eqref{efflerch:eq:mesh}, including strict simplicity of the limit.
\end{proof}

\begin{lemma}[The inherited global zero bound]
\label{efflerch:lem:variation}
For $n\geq0$, $k\geq1$, and $0\leq\rho\leq1$,
$F_{n,k}^\rho$ has at most $n$ positive zeros, counted with multiplicity.
\end{lemma}
\begin{proof}
We recall the short variation argument so that the completeness claim
does not depend on a root scan.  If a real kernel $g$ has
$N$ sign changes, its Laplace transform has at most $N$ positive zeros,
provided the exponentially weighted polynomial moments exist.
For $N=0$ the transform has a strict constant sign.  For $N>0$, choose
a sign-change point $t_0$.  Multiplication by $t_0-t$ removes that
change.  The new transform is
\[
 (\partial_a+t_0)G(a)
   =e^{-t_0a}\frac{d}{da}\bigl(e^{t_0a}G(a)\bigr).
\]
By induction it has at most $N-1$ zeros.  Any finite collection of
$M$ zeros of $G$, counted with multiplicity, forces at least $M-1$
zeros of the displayed derivative by Rolle's theorem with multiplicities.
Thus $M\leq N$.  Induction also precludes an identically zero transform.

In \eqref{efflerch:eq:F}, the factor $t^kW_\rho(t)$ is positive and
$Q_n(\log t)$ has $n$ simple sign changes by
Proposition~\ref{efflerch:prop:mesh}, with the cases $n=0,1$ immediate.
The bound $W_\rho(t)\leq W_1(t)\leq1+t^{-1}$ gives all required
integrability near zero, and $e^{-at}$ controls infinity.
\end{proof}

\subsection{A scale-free positive-weight estimate}

The decisive improvement over a compact-scale Taylor estimate is the
following uniform control of the weight on the logarithmic coordinate.

\begin{lemma}[Logarithmic variation of the Lerch weight]
\label{efflerch:lem:weight}
For $x,y>0$ and $0\leq\rho\leq1$,
\begin{equation}\label{efflerch:eq:weight}
 \left|\log\frac{W_\rho(xy)}{W_\rho(x)}\right|\leq|\log y|.
\end{equation}
\end{lemma}
\begin{proof}
Direct differentiation gives
\[
 \frac{d\log W_\rho(x)}{d\log x}
  =-\frac{x\rho}{e^x-\rho}.
\]
Its magnitude is at most one: since $e^x\geq1+x$ and $0\leq\rho\leq1$,
$e^x-\rho\geq1+x-\rho\geq\rho x$.
Integrate this derivative between $\log x$ and $\log x+\log y$.
The proof is valid at both deformation endpoints.
\end{proof}

Let $Y_k$ have gamma density
\[
 \frac{k^{k+1}}{k!}y^k e^{-ky},\qquad y>0.
\]
This is shape $k+1$ and rate $k$.  Keeping the shape shift is important:
for example $\mathbb E Y_k^{-1}=1$ exactly.

\begin{lemma}[An elementary exponential log-moment bound]
\label{efflerch:lem:moments}
For positive integers $m,k$ with $k\geq16m^2$,
\begin{equation}\label{efflerch:eq:moments}
 \mathbb E\bigl(e^{m|\log Y_k|}-1\bigr)
 \leq\left[2\bigl(e^{4m^2/k}-1\bigr)\right]^{1/2}
 \leq\sqrt{2/3}<1.
\end{equation}
\end{lemma}
\begin{proof}
For a positive number $y$, according to whether $y\geq1$ or $y<1$,
the square of $e^{m|\log y|}-1$ is one of
$(y^m-1)^2$ and $(y^{-m}-1)^2$.  Cauchy--Schwarz therefore bounds
the square of the left side by
\[
 \mathbb E Y_k^{2m}+\mathbb E Y_k^{-2m}
  -2\mathbb E Y_k^m-2\mathbb E Y_k^{-m}+2.
\]
The integer gamma moments are
\[
 \mathbb E Y_k^r=\prod_{j=1}^r(1+j/k),\qquad
 \mathbb E Y_k^{-r}=\prod_{j=0}^{r-1}(1-j/k)^{-1}\quad(r\leq k).
\]
Both are at least one.  Since $k\geq4m$,
\begin{align*}
 \log\mathbb E Y_k^{2m}
  &\leq\frac{m(2m+1)}k\leq\frac{3m^2}k,\\
 \log\mathbb E Y_k^{-2m}
  &\leq\frac2k\sum_{j=0}^{2m-1}j
   \leq\frac{4m^2}k.
\end{align*}
Here $-\log(1-u)\leq2u$ for $0\leq u\leq1/2$ supplies the second
line.  These estimates prove the first inequality in
\eqref{efflerch:eq:moments}.  Finally $4m^2/k\leq1/4$ and
$e^{1/4}\leq(1-1/4)^{-1}=4/3$ give the displayed rational bound.
\end{proof}

\begin{proposition}[Uniform relative sign control away from Appell roots]
\label{efflerch:prop:sign}
Fix $n,k\geq1$, a real $s$ at distance at least $d>0$ from every
root of $Q_n$, and put $m=\lceil1+n/d\rceil$.
If $k\geq16m^2$, then, simultaneously for $0\leq\rho\leq1$,
\begin{equation}\label{efflerch:eq:relative}
 \left|
  \frac{(ke^{-s})^{k+1}F_{n,k}^\rho(ke^{-s})}
       {k! W_\rho(e^s)Q_n(s)}-1
 \right|\leq\sqrt{2/3}<1.
\end{equation}
In particular, $F_{n,k}^\rho(ke^{-s})$ and $Q_n(s)$ have the same sign.
\end{proposition}
\begin{proof}
The substitution $t=e^s y$ in \eqref{efflerch:eq:F} gives the exact
normalized expression
\[
 \frac{(ke^{-s})^{k+1}F_{n,k}^\rho(ke^{-s})}
      {k!W_\rho(e^s)Q_n(s)}
 =\mathbb E\left[
  \frac{W_\rho(e^sY_k)}{W_\rho(e^s)}
  \frac{Q_n(s+\log Y_k)}{Q_n(s)}\right].
\]
For $u\in\mathbb R$, the product over roots gives
\[
 \left|\frac{Q_n(s+u)}{Q_n(s)}-1\right|
  \leq\prod_{j=1}^n\left(1+\frac{|u|}{|s-x_{n,j}|}\right)-1
  \leq e^{n|u|/d}-1.
\]
Put $R=W_\rho(e^{s+u})/W_\rho(e^s)$.  Lemma~\ref{efflerch:lem:weight}
implies $R\leq e^{|u|}$ and $|R-1|\leq e^{|u|}-1$.
Consequently
\[
 \left|R\frac{Q_n(s+u)}{Q_n(s)}-1\right|
  \leq e^{(1+n/d)|u|}-1\leq e^{m|u|}-1.
\]
Take expectations and apply Lemma~\ref{efflerch:lem:moments}.
This also supplies all needed absolute integrability.
\end{proof}

\subsection{Completion of the zero theorem and the joint limit}

\begin{proof}[Proof of Theorem~\ref{efflerch:thm:explicit}]
Put $\eta=8n/\sqrt{k}$.  By \eqref{efflerch:eq:Kn},
$\eta\leq\delta_n/3$ and $k\geq256$.
The intervals $[x_{n,j}-\eta,x_{n,j}+\eta]$ are disjoint.
Each endpoint is at distance at least $\eta$ from every Appell root.
For $d=\eta$ the integer in Proposition~\ref{efflerch:prop:sign} satisfies
\[
 m=\left\lceil1+\frac{\sqrt{k}}8\right\rceil
  \leq2+\frac{\sqrt{k}}8\leq\frac{\sqrt{k}}4,
\]
so $k\geq16m^2$.  That proposition transfers the opposite Appell
signs at the two endpoints to the normalized Lerch family, uniformly
in $\rho$.  The decreasing change of coordinate $a=ke^{-s}$ gives
one positive zero in each of the corresponding $n$ disjoint intervals.

The global zero bound counts multiplicity and allows at most $n$.
It therefore excludes every additional positive zero and every
multiple zero.  The $j$th interval consequently contains precisely
$a_{n,j,k}^\rho$ in the stated ordering.  Its endpoints yield
\eqref{efflerch:eq:location}.
\end{proof}

\begin{corollary}[A growing Laurent index]
\label{efflerch:cor:growing}
Fix $0<\alpha<1/2$.  For all sufficiently large integers $k$,
simultaneously for
\[
 2\leq n\leq\alpha\frac{\log k}{\log\log k},\qquad
 1\leq j\leq n,\qquad 0\leq\rho\leq1,
\]
the family has exactly $n$ simple positive zeros and
\begin{equation}\label{efflerch:eq:growing}
 \frac{a_{n,j,k}^\rho}{k e^{-x_{n,j}}}=1+o(1).
\end{equation}
More precisely, the logarithmic error is bounded by $8n/\sqrt{k}$
throughout this range.
\end{corollary}
\begin{proof}
Equation~\eqref{efflerch:eq:Kn} gives
$\log K_n=2n\log n+O(n)$.  At the largest permitted $n$, this is
at most $(2\alpha+o(1))\log k<\log k$; the same bound holds for
every smaller $n$, with finitely many bounded indices absorbed by
increasing the lower threshold for $k$.  Thus $k\geq K_n$ throughout
the range.  Theorem~\ref{efflerch:thm:explicit} applies, and
$8n/\sqrt{k}\to0$ uniformly there.  Exponentiation proves
\eqref{efflerch:eq:growing}.
\end{proof}

\begin{remark}[A replaceable separation input]
The argument actually proves the stronger reusable statement: if any
$\delta>0$ is a proved lower bound for $\operatorname{mesh}(Q_n)$, then
the same location bound holds once
\[
 k\geq\max\{256,(24n/\delta)^2\}.
\]
The explicitly constructed $\delta_n$ makes this unconditional for every
index.  Sharper analytic estimates or certified Appell root brackets can
be inserted without changing the positive-kernel proof.  This is a
natural route to much smaller effective thresholds.
\end{remark}

\subsection{Concrete sufficient thresholds from rational root isolation}

The proof of Proposition~\ref{efflerch:prop:mesh} gives the more precise
inequality $\operatorname{mesh}(Q_n)\geq\operatorname{mesh}(B_n)$.
The latter is a rational-polynomial problem.  The accompanying
\texttt{lerch\_verify\_effective.py} constructs $B_n$ in exact fraction
arithmetic, removes its simple zero at the origin, and isolates every
remaining root by a Sturm sequence.  If $[\ell_i,u_i]$ are the ordered
isolating intervals, then
$\min_i(\ell_{i+1}-u_i)$ is a rational lower mesh bound.
The exact root count certifies completeness of this finite list.

The table records a subset of the replayed cases $2\leq n\leq12$.
Its last column is exactly
$\left\lceil\max\{256,(24n/\underline\delta)^2\}\right\rceil$.
The executable checks both the upper-ceiling inequality and its minimality
as an integer ceiling by rational comparisons.
\begin{center}
\begin{tabular}{rcc}
\hline
$n$ & Certified $\underline\delta\leq\operatorname{mesh}(B_n)$
 & Sufficient $k$ for every $\rho\in[0,1]$\\
\hline
2 & $16383/8192$ & $577$\\
3 & $3333/4096$ & $7830$\\
4 & $29421/65536$ & $45729$\\
5 & $4683/16384$ & $176261$\\
6 & $6495/32768$ & $527797$\\
7 & $9557/65536$ & $1327197$\\
8 & $3663/32768$ & $2950044$\\
\hline
\end{tabular}
\end{center}
These are sufficient thresholds obtained from a common quantitative
argument; they are not sharp classifications.  The known sharp uniform
threshold at $n=2$ remains $k=1$.  The same replay verifies six exact
gamma-moment controls at $k=16m^2$, for $m=1,2,4,8,16,32$, and rejects
a deliberately out-of-hypothesis control at $m=2,k=4$.
These finite checks audit the implementation.  The all-index theorem
and its growing-index corollary rest on the preceding proofs.

\subsection{The terminal branch at arbitrary index}

One further general consequence does not require concentration.  It
extends the known index-two terminal-branch motion law to any index and
derivative order whenever the last positive zero is simple.

\begin{proposition}[Strict leftward motion of a simple terminal zero]
\label{efflerch:prop:terminal}
Fix $n,k\geq1$.  On an interval $I\subset[0,1)$ suppose the largest
positive zero $r(\rho)$ of $F_{n,k}^\rho$ exists, is simple, and is
followed by no further positive zero.  Then $r'(\rho)<0$ on the interior
of $I$.  At $\rho=0$ the same conclusion holds for a right derivative
if the simple terminal branch persists there.  In particular, for
$n\geq2$ and $k\geq K_n$ the last branch is strictly decreasing on $[0,1]$.
\end{proposition}
\begin{proof}
For $0\leq\rho<1$, absolute convergence gives the resolvent identity
\begin{equation}\label{efflerch:eq:resolvent}
 \partial_\rho F_{n,k}^\rho(a)
  =\sum_{h\geq1}\rho^{h-1}F_{n,k}^\rho(a+h).
\end{equation}
Indeed, write $F_{n,k}^\rho(a)=\sum_{m\geq0}\rho^m f(a+m)$,
where $f=F_{n,k}^0$, and sum over $h+m$; the coefficient of
$\rho^{q-1}f(a+q)$ is $q$.  Every summand on the right of
\eqref{efflerch:eq:resolvent}, evaluated at $a=r(\rho)$, has the
same strict sign as $F_{n,k}^\rho$ just to the right of this last
simple zero.  That is also the sign of $\partial_aF_{n,k}^\rho(r(\rho))$.
Implicit differentiation yields $r'=-F_\rho/F_a<0$.

In the saturated regime all branches are simple for every
$\rho\in[0,1]$.  The uniform compact brackets supplied by
\eqref{efflerch:eq:location}, continuity of the integral in $\rho$,
and uniqueness inside each bracket give continuity at both endpoints.
Strict decrease on the interior therefore extends to the closed interval.
No finite $\rho$-derivative at $\rho=1$ is asserted here.
\end{proof}

